\documentclass[a4paper,12pt,twoside,final]{scrbook}
\usepackage[utf8]{inputenc}
\usepackage[spanish,es-nodecimaldot]{babel} % International, rec. by RAE: http://www.tex-tipografia.com/marca_decimal.html

\usepackage{booktabs}
\usepackage[hidelinks]{hyperref}

\usepackage{tgpagella} % text only
\usepackage{mathpazo}  % math & text




\usepackage{setspace}
\spacing{1.2}

% COLOR
\usepackage[usenames,dvipsnames,svgnames,table]{xcolor}
% Colores para estilo Proyecto Docente (tonos naranjas)
\definecolor{lightback}{HTML}{F4E0BF}
     \definecolor{back}{HTML}{F3C591}
\definecolor{lightline}{HTML}{FCAF5F}
     \definecolor{line}{HTML}{ED7900}
% Colores para portada
\definecolor{epsc:oscuro}{HTML}{280091}
 \definecolor{epsc:medio}{HTML}{4C5CC5}
 \definecolor{epsc:claro}{HTML}{3FCFD5}
 \definecolor{epsc:verde}{HTML}{00B299}

\usepackage{tikz} % used in cover to place images
\usepackage[simplified]{pgf-umlcd}
\usepackage{pgf-umlsd}

\usepackage{datetime} % allow formal date format
% "Month, YEAR" date format, in spanish with the month uppercased not interfering other dates
\newcommand\Monthname[1][EMPTY]{
  \ifnum1=#1Enero\else
  \ifnum2=#1Febrero\else
  \ifnum3=#1Marzo\else
  \ifnum4=#1Abril\else
  \ifnum5=#1Mayo\else
  \ifnum6=#1Junio\else
  \ifnum7=#1Julio\else
  \ifnum8=#1Agosto\else
  \ifnum9=#1Septiembre\else
  \ifnum10=#1Octubre\else
  \ifnum11=#1Noviembre\else
  \ifnum12=#1Diciembre\else
  \fi\fi\fi\fi\fi\fi\fi\fi\fi\fi\fi\fi
}
\newdateformat{monthyeardate}{%
  \Monthname[\THEMONTH], \THEYEAR}


% FONTs
\usepackage[T1]{fontenc}
\usepackage{textcomp}           % Needed for new symbols like € ?

\usepackage[scaled]{berasans} % Font for the cover similar to Vera 33
%\renewcommand*\familydefault{\sfdefault}  %% To use as the base font of the document is to be sans serif


% PAGE STYLE
\usepackage[twoside,bindingoffset=0mm,headheight=16pt,margin=25mm]{geometry} %,verbose,showframe
%\usepackage{fancyhdr} % Encabezados
%\pagestyle{fancy}
% XXX Evitar el binding en la portada pero no en el resto del documento

% TABLES AND FIGURES
\usepackage{graphicx}
\graphicspath{ {img/} }

% SUBRALLADOS
\usepackage[normalem]{ulem}



\begin{document}
%\frontmatter

%------------- Cover --------------
\thispagestyle{empty}

% Backgroud images
\begin{tikzpicture}[remember picture, overlay]
  % Top
  \node [anchor=north east, inner sep=0pt]  at (current page.north east)
     {\includegraphics[height=6cm]{topRightCorner.pdf}};
  % Bottom
  \node [anchor=south west, inner sep=0pt]  at (current page.south west)
     {\includegraphics[height=6cm]{bottomLeftCorner.pdf}};
  \node (uco) [anchor=south east, inner sep=0pt, xshift=-10mm, yshift=10mm]  at (current page.south east)
        {\includegraphics[height=2cm]{uco_debajo.pdf}};
  \node [anchor=south east, inner sep=0pt, xshift=-10mm]  at (uco.south west)
% Uncomment the chosen logo and comment the others:
        {\includegraphics[height=2cm]{emblema-ing-informatica.pdf}};
%        {\includegraphics[height=2cm]{emblema-ing-industrial.pdf}};
%        {\includegraphics[height=2cm]{emblema-ing-tec-industrial.pdf}};
\end{tikzpicture}


\begin{center}
\fontfamily{\sfdefault}\selectfont
\vspace*{2cm}

\vfill
\vfill
\includegraphics[width=12.5cm]{LogotipoEPSC.pdf}
\vfill
\vfill

\large\textbf{\color{epsc:medio}
  TRABAJO FIN DE GRADO EN INGENIERÍA INFORMÁTICA
}
\vfill

\Large\textbf{\color{epsc:verde}
  MEMORIA TÉCNICA
}
\vfill
\vfill

\Huge\textbf{\color{epsc:oscuro}
  Plataforma web de gestión de comunicados
}
\vfill
\vfill

\large{\color{epsc:oscuro}Autor}\\
\textbf{\color{epsc:medio}{ Eduardo Arroyo Ramírez }}
\vfill

\large{\color{epsc:oscuro} Director }\\
\textbf{\color{epsc:medio} Domingo Ortiz Boyer}
\vfill



\textbf{\color{epsc:verde} \monthyeardate\today}
\vfill
\vfill
\vspace{2.7cm}
\end{center}

%-------------------------------------------------------------------------------
\clearpage

\thispagestyle{empty}
\pagecolor{white}

\cleardoublepage
%-------------------------------------------------------------------------------
\frontmatter
\tableofcontents
\listoffigures
\listoftables

\cleardoublepage
\mainmatter
\chapter{Introducción}
\section{Definición del problema}
La \emph{Comunicación Corporativa} se define como <<la totalidad de los recursos de comunicación de los que dispone una organización para llegar efectivamente a sus públicos>> \cite{riel_2001}. Dentro de este concepto, se distingue entre la \emph{Comunicación Corporativa Interna} y la \emph{Comunicación Corporativa Externa} siendo la primera <<todo lo relativo a la conexión requerida entre los miembros de una determinada estructura para acometer una metas comunes>> y la segunda, <<la vinculación de la organización con el entorno en el que desarrolla sus actividades, con el fin de alcanzar un determinado nivel de rentabilidad económica y social>> \cite{castro_2007}. Desde el ámbito empresarial y administrativo se da cada vez mayor importancia a la comunicación ya que ésta juega un papel fundamental en la consecución de sus objetivos.

Este tipo de comunicación ha encontrado en Internet uno de sus canales principales y tanto empresas como instituciones ejercen hoy en día su presencia online a través de webs, blogs, redes sociales o incluso sus propias aplicaciones para móviles, que han proliferado enormemente durante los últimos años \cite{playstore} \cite{appstore}. Todas estas herramientas permiten no sólo la publicación de comunicados sino también la planificación temporal, seguimiento y medida del impacto.

Aunque las redes sociales son uno de los medios preferidos por las empresas \cite{linkedin} y administraciones \cite{grande2015casos} como canal de comunicación debido a la gran acogida que tienen entre la población, las aplicaciones móviles son una alternativa muy a tener en cuenta ya que cubren la muy frecuente necesidad de incluir otros servicios específicos de la organización como por ejemplo:
\begin{itemize}
    \item Sistema de reservas (mesas en un restaurante, salas en un gimnasio\dots).
    \item Encuestas (toma de decisiones en una comunidad de vecinos, sondeo de la clientela de un bar para organizar un concierto en directo\dots)
    \item Sistema de geolocalización en tiempo real de un desfile (¿por dónde está pasando ahora mismo la cabalgata de reyes?).
\end{itemize}
Además, las aplicaciones salvan algunos de los inconvenientes de las redes sociales como la cesión de datos, la alta exposición a la crítica o la obligación de adaptarse al diseño y las reglas de negocio impuestas por un tercero.

\section{Antecedentes}

\cleardoublepage
\chapter{Objetivos}
El objetivo general de este trabajo es la creación de una plataforma software que permita ofrecer en poco tiempo y a coste reducido aplicaciones móviles personalizadas asociadas a un backend que permita a diferentes organizaciones realizar la gestión básica y seguimiento del impacto de sus comunicaciones corporativas.

Como objetivo secundario, la arquitectura de la plataforma facilitará la incorporación de módulos para los servicios personalizados que cada organización pueda solicitar. Estos módulos podrán estar hechos a medida para una organización o bien estar previamente desarrollados de forma que se activen o desactiven dependiendo de si una organización está interesada en la funcionalidad que aportan.

Para la consecución de dichos objetivos, se establecen los siguientes objetivos específicos:

\section{Plataforma de publicación y seguimiento de comunicados}
Es necesario dotar al sistema de una herramienta que permita a los editores elaborar, publicar, programar (publicación diferida) y consultar las métricas de impacto de los comunicados. Para ello se desarrollará una aplicación web con acceso individualizado que cubrirá las siguientes necesidades según el tipo de usuario:
\begin{itemize}
    \item Lectores (acceso público)
    \begin{itemize}
        \item Consulta pública de comunicados de la organización.
        \item Consulta de otros datos de interés de la organización (teléfonos, direcciones, etc.).
    \end{itemize}
    \item Editores (acceso restringido por organización)
    \begin{itemize}
        \item Administración de comunicados.
        \item Consulta de métricas de impacto de comunicaciones.
        \item Administración de tipos de registro independientes como categorías de comunicaciones, teléfonos, direcciones, etc. de la organización.
    \end{itemize}
    \item Administradores (responsables del mantenimiento del servicio)
    \begin{itemize}
        \item Administración de organizaciones.
        \item Administración de funcionalidades personalizadas disponibles para cada organización.
        \item Administración de usuarios.
        \item Configuración del servicio.
    \end{itemize}
\end{itemize}

\section{Servicio WEB de consulta de comunicados}
Para brindar acceso a los comunicados a través de aplicaciones propias o de terceros, es necesario dotar a la plataforma de un servicio WEB con una API pública que ofrezca los métodos necesarios para realizar dicha labor.
\begin{itemize}
    \item Lectores (acceso público)
    \begin{itemize}
        \item Consulta pública de comunicados de la organización.
        \item Consulta de otros datos de interés de la organización (teléfonos, direcciones, etc.).
    \end{itemize}
\end{itemize}

\section{Aplicación móvil tipo}
Se desarrollará una aplicación móvil que permita consultar los comunicados de una organización desde su teléfono móvil, consumiendo el servicio web de la plataforma mediante solicitudes HTTP. Esta aplicación será un prototipo precursor de las aplicaciones para clientes y demostrará las capacidades del sistema.
\cleardoublepage
\chapter{Recursos}
En esta sección se clasificarán y detallarán los recursos necesarios para la realización del trabajo.

\section{Humanos}
En este apartado se identifica a las personas involucradas en el trabajo.
\begin{itemize}
    \item Prof. D. \emph{Domingo Ortiz Boyer} profesor del Departamento de Informática y Análisis Numérico de la Escuela Politécnica Superior de la Universidad de Córdoba, como tutor del proyecto.
    \item D. \emph{Eduardo Arroyo Ramírez}, encargado de realizar dicho proyecto siguiendo las directrices definidas por el tutor del mismo.
\end{itemize}

\section{Hardware}
Los recursos hardware son las herramientas físicas empleadas para el desarrollo y la puesta en marcha del proyecto.
\begin{itemize}
    \item Ordenador PC Compatible con procesador Intel i7 y 8GB de memoria RAM.
    \item Teléfono móvil Android.
    \item Conexión a Internet de banda ancha.
\end{itemize}

\section{Software}
Bajo esta denominación se engloban las herramientas software como el sistema operativo, los entornos de programación, bibliotecas, herramientas de documentación, etc.
\begin{itemize}
    \item S. O. Windows 10 Home Premium
    \item Microsoft Visual Studio Community 2019
    \item Telerik UI para ASP.NET MVC
    \item Microsoft SQL Server Express 2016
    \item Titanium Studio
    \item Internet Information Server 7
    \item Navegadores:
    \begin{itemize}
        \item Google Chrome
        \item Mozilla Firefox
        \item Microsoft Edge
    \end{itemize}
    \item Microsoft Visual Studido Code
    \item Paquete \LaTeX de edición de textos
\end{itemize}
\cleardoublepage
\chapter{Análisis de requisitos}
Este apartado se centra en la investigación del problema y los requisitos. Aquí se definirán los conceptos en el dominio del problema.

\section{Modelo de casos de uso}

¿Qué es un modelo de casos de uso?
Formato utilizado
Listado de casos de uso formato breve con identificador único
Con el fin de identificar los actores, conceptos y actividades del sistema se elaborará un \emph{modelo de casos de uso} \cite{Jacobson1999} \cite{Larman2004}

\begin{itemize}
    \item Publicar un comunicado
    \item Programar un comunicado
    \item Consultar law lista de comunicados
    \item Consultar la lista de comunicados por categoría
    \item Consultar un detalle de comunicado
    \item 
\end{itemize}


\section{Modelo de dominio}

\begin{tikzpicture}[show background grid]
    \begin{class}[text width=6cm]{Class}{0 ,0}
        \attribute{+Public }
        \attribute{\#Protected }
        \attribute{-Private }
        \attribute{$\sim$ Package }
    \end{class}
    \begin{class}[text width=7cm]{BankAccount}{7,0}
        \attribute{+owner:String}
        \attribute{+balance:Dollar s}
        \operation{+deposit(amount:Dollars)}
        \operation{+withdrawal(amount:Dollars)}
        \operation{\#updateBalance(newBalance:Dollars)}
    \end{class}
\end{tikzpicture}

\begin{sequencediagram}
    \newthread{t}{:Thread}
    \newinst{a}{:A}
    \newinst{b}{:B}
    \begin{call}{t}{funcA()}{a}{return}
        \begin{call}{a}{funcA()}{b}{return}
        \end{call}
    \end{call}
\end{sequencediagram}
\cleardoublepage
\include{capitulos/050-diseño}
\cleardoublepage
\chapter{Implementación}

La parte WEB de este proyecto  --incluído el servicio para las apps-- estará construida con .NET MVC, un framework de Microsoft para el desarrollo de aplicaciones WEB que implementa el patrón de diseño modelo-vista-controlador. .NET MVC es software libre y se distribuye bajo licencia Apache 2.0.
\cleardoublepage
\chapter{Pruebas}
\cleardoublepage
\chapter{Futuras mejoras}

\section{Capa de servicios}
Exponer clases POCO a modo de vistas en lugar de exponer las entidades. Es decir, PersonasServicio.BuscarPersona() debe devolver PersonaModel y no PersonaEntity.

\section{Aplicación web}

\subsection{Backend}
Inyección de dependencias
Prescindir de OWIN

\subsection{Frontend}
Quitar Telerik
\cleardoublepage
\include{capitulos/100-conclusiones}
\cleardoublepage

% BIBLIOGRAFÍA -----------------------------------------------------------------
\bibliographystyle{ieeetr}
\renewcommand{\refname}{Bibliografía}
\bibliography{00-memoria-tecnica}
\addcontentsline{toc}{part}{Bibliografía}
\clearpage
% ------------------------------------------------------------------------------

\end{document}